\section{Database Management Systems}
\label{sec:qna-dbms}

\subsection*{Foundations}

\question{What is a DBMS and what problems does it solve over flat files?}
\answer{A DBMS manages structured data with querying, concurrency control,
integrity constraints, transactions, and recovery. Over flat files it adds
data independence, controlled redundancy, multi-user safety (ACID), and a
declarative query language, removing the need for each application to
reimplement these.}

\question{Difference between primary key, candidate key, and foreign key?}
\answer{A candidate key is any minimal set of columns uniquely identifying
a row. The primary key is the chosen candidate key (unique, non-null). A
foreign key is a column referencing another table's primary key, enforcing
referential integrity between tables.}

\question{What is normalization? Explain 1NF, 2NF, 3NF.}
\answer{Normalization organises columns to reduce redundancy and anomalies.
\textbf{1NF}: atomic values, no repeating groups. \textbf{2NF}: 1NF plus no
partial dependency of a non-key column on part of a composite key.
\textbf{3NF}: 2NF plus no transitive dependency (non-key columns depend
only on the key). BCNF is a stricter 3NF.}

\question{What is denormalization and when is it used?}
\answer{Deliberately introducing redundancy (e.g.\ storing a computed or
duplicated column) to speed up reads by avoiding joins. Used in
read-heavy/analytics systems where query performance outweighs the cost of
maintaining consistency on writes.}

\subsection*{Transactions and concurrency}

\question{What are ACID properties?}
\answer{\textbf{Atomicity} (all-or-nothing), \textbf{Consistency} (a
transaction moves the DB from one valid state to another),
\textbf{Isolation} (concurrent transactions don't interfere as if serial),
\textbf{Durability} (committed changes survive crashes). Together they make
transactions reliable.}

\question{What are the SQL isolation levels and what anomalies do they
prevent?}
\answer{From weakest to strongest: \emph{Read Uncommitted} (allows dirty
reads), \emph{Read Committed} (prevents dirty reads), \emph{Repeatable
Read} (also prevents non-repeatable reads), \emph{Serializable} (also
prevents phantom reads). Higher isolation means more correctness but less
concurrency.}

\question{What is a deadlock and how do databases handle it?}
\answer{Two or more transactions each wait for locks the others hold,
forming a cycle so none proceed. DBMSs detect it (wait-for graph) and abort
a victim, or prevent it via ordering/timeouts. It is the transactional
analogue of an OS deadlock.}

\question{Difference between optimistic and pessimistic locking?}
\answer{Pessimistic locking locks data before use, assuming conflicts are
likely -- safe but reduces concurrency. Optimistic locking assumes
conflicts are rare: it checks at commit (via a version/timestamp) whether
data changed, and retries if so. Optimistic suits low-contention,
read-heavy workloads.}

\subsection*{Indexing and querying}

\question{What is an index and what is the trade-off?}
\answer{An index is an auxiliary structure (usually a B+ tree) that speeds
up lookups and range queries on indexed columns, turning a full scan into a
logarithmic search. The trade-off: extra storage and slower
writes/updates, since indexes must be maintained.}

\question{Clustered vs non-clustered index?}
\answer{A clustered index determines the physical order of rows -- there is
at most one per table, and the leaf level \emph{is} the data. A
non-clustered index is separate and stores pointers (or the clustered key)
to the rows; a table can have many. Clustered lookups avoid an extra
indirection.}

\question{What is the difference between WHERE and HAVING?}
\answer{\texttt{WHERE} filters rows \emph{before} grouping; \texttt{HAVING}
filters groups \emph{after} \texttt{GROUP BY} aggregation. You use
\texttt{HAVING} to filter on aggregate results like
\texttt{COUNT(*) > 5}.}

\question{Difference between DELETE, TRUNCATE, and DROP?}
\answer{\texttt{DELETE} removes rows (optionally with a \texttt{WHERE}),
is logged, and can be rolled back. \texttt{TRUNCATE} removes all rows
quickly with minimal logging and resets identity, but is harder to roll
back. \texttt{DROP} removes the entire table structure and data.}

\subsection*{Scale}

\question{SQL vs NoSQL -- when would you choose each?}
\answer{SQL (relational) suits structured data, complex joins, and strong
consistency/ACID -- e.g.\ finance. NoSQL (document, key-value, column,
graph) suits flexible schemas, huge scale, and high write throughput with
eventual consistency -- e.g.\ user feeds, telemetry. Choose by consistency
needs, schema stability, and access patterns.}

\question{What is the CAP theorem?}
\answer{In a distributed store you cannot simultaneously guarantee all of
Consistency, Availability, and Partition tolerance; under a network
partition you must trade consistency for availability or vice versa.
Real systems (partition tolerance being mandatory) are effectively CP or
AP.}
